\chapter{Podsumowanie}
\label{cha:podsumowanie}

Celem projektu było stworzenie aplikacji wspierającej nauczyciela i ucznia w pełnym cyklu pracy
z testami - od ich ułożenia, przez udostępnienie i rozwiązanie, aż po ocenę i analizę wyników. Cel
ten został zrealizowany: powstała spójna, w pełni funkcjonalna aplikacja desktopowa napisana
w języku C\# na platformie .NET, z interfejsem graficznym opartym na technologii WPF i z relacyjną
bazą danych PostgreSQL obsługiwaną przez Entity Framework Core.

\section*{Zrealizowane funkcjonalności}

W ramach projektu udało się zrealizować wszystkie założone funkcje:

\begin{itemize}
  \item \textbf{Konta i role} - rejestracja (z wyborem roli) i logowanie użytkowników
        z rozróżnieniem ról nauczyciela i ucznia oraz odrębnymi panelami dla każdej z nich; każdy
        nauczyciel zarządza wyłącznie własnymi testami.
  \item \textbf{Grupy} - tworzenie grup i zarządzanie ich składem, z możliwością przynależności
        ucznia do wielu grup.
  \item \textbf{Testy} - tworzenie, edycja, duplikowanie i usuwanie testów złożonych z pytań
        jednokrotnego wyboru, wielokrotnego wyboru oraz pytań otwartych, wraz z walidacją
        kompletności testu oraz sterowaniem widocznością (test aktywny/nieaktywny).
  \item \textbf{Rozwiązywanie testów} - rozwiązywanie testów w ograniczonym czasie, z obsługą
        wielu podejść i wyborem najlepszego wyniku.
  \item \textbf{Ocenianie} - automatyczne ocenianie pytań zamkniętych (z punktami częściowymi)
        oraz ręczne ocenianie i edycja oceny dowolnego podejścia (z możliwością nadpisania punktów
        każdego pytania) z automatycznym przeliczaniem wyniku.
  \item \textbf{Bezpieczeństwo i integralność} - haszowanie haseł, unikalność loginów, wykrywanie
        prób oszustwa oraz spójność danych zapewniona przez klucze obce i migracje.
  \item \textbf{Analiza wyników} - statystyki testu z prezentacją wszystkich podejść uczniów oraz
        eksport wyników do pliku CSV.
\end{itemize}

W trakcie tworzenia projektu zastosowano zasady programowania obiektowego: dziedziczenie
(odwzorowane w bazie danych strategią Table-per-Hierarchy), enkapsulację oraz podział kodu na
niezależne, logiczne warstwy. Dzięki temu struktura aplikacji jest czytelna i łatwa w utrzymaniu.

\section*{Możliwości dalszego rozwoju}

Choć aplikacja jest w pełni funkcjonalna, można ją w przyszłości rozbudować, na przykład o:

\begin{itemize}
  \item losowanie kolejności pytań i odpowiedzi w celu dodatkowego utrudnienia ściągania,
  \item podgląd historii własnych podejść po stronie ucznia (obecnie uczeń widzi swój najlepszy
        wynik, a pełną historię podejść ma nauczyciel w statystykach),
  \item terminy dostępności testów (przedział dat, w którym test można rozwiązać),
  \item rozbudowane raporty i wykresy postępów uczniów w czasie.
\end{itemize}

Dodanie tych funkcji zwiększyłoby użyteczność systemu, jednak w obecnej formie projekt w pełni
spełnia założone wymagania i stanowi solidną podstawę do dalszej rozbudowy.
